\chapter*{How to Read This Book}
\addcontentsline{toc}{chapter}{How to Read This Book}
\markboth{HOW TO READ THIS BOOK}{HOW TO READ THIS BOOK}

\begin{flushright}
\itshape Lasciate ogne fretta, voi ch'intrate.\\[2pt]
\upshape\small not quite Dante, \emph{Inferno}, canto III
\end{flushright}

Over the gate of Hell, Dante wrote: abandon all hope, ye who enter here. This book asks
for less: abandon all haste.

This is a dense book. It is packed with names, dates, formulas and numbers, and almost
every page leans on something said before it. It assumes you will stop, go back a few
pages, check a number, and come back. It was written for curious people: the ones who
want to know why a thing is the way it is, who would rather spend an evening on one page
than skim twenty, and who are not satisfied until the idea and the history behind it
both make sense.

One question runs through the whole book: how far are we from a quantum computer
that solves a useful problem, and which problem would that be? Answering it means
building the language first and then going down a stack, one layer at a time. The
book has one part for the language and one for each layer, and each part answers one
question the main one forces on you.

The first part is the foundations, and the question is what a quantum computer is.
It starts where the field started, with Richard Feynman asking in 1981 why we
simulate nature with machines that do not work the way nature does. Then it builds
the rest: what a computer can do at all, what a qubit is, where quantum mechanics
came from, and the linear algebra underneath it.

The second part is about the problems: which real problems are waiting for such a
machine, and how each one turns into something it can act on. There are several, and
they are not equally promising.

The third part is about the algorithm: what the circuit extracts from the problem,
and at what cost.

The fourth part is about hardware: which technology delivers physical qubits, and
with how much error.

The fifth part is about error correction: how many of those physical qubits it takes
to make one qubit you can trust.

The sixth part is about the market: who sells each layer, and what you can actually
use today.

The parts are in the order in which the questions depend on each other, but the
chapters are written in the order the study reaches each subject. A part can grow
while the one before it is still thin, and when a chapter needs something from
another part, it says where to find it.

\section*{Every idea has a parent}

No idea in this book showed up because it was elegant. Max Planck put energy in
packets in 1900 because nothing else fitted the light coming out of a hot body, and
he called it an act of desperation. Kenneth Wilson put quarks on a lattice in 1974
because the strong force was too strong for the approximations that had worked for
electrodynamics. Peter Shor found his factoring algorithm in 1994 by thinking about a
paper by Daniel Simon that he had helped reject. Shor's nine-qubit code exists because
Rolf Landauer said quantum computers could never correct their errors. The transmon
was designed in 2007 because the superconducting qubits before it were ruined by stray
electric charges. Each idea fixed something its predecessor got wrong and brought
costs of its own. If you know the lineage, you know when the idea applies. If you only
know the rule, you will apply it where it does not belong.

\section*{Five kinds of boxes}

Stories, concepts, questions, reflections and words are kept apart, so you
always know which one you are reading.

\begin{history}{year \textperiodcentered{} place}
An amber box tells a story: who did it, when, and what forced the idea into
existence. The people, the dates and the problems in it are real. No invented
incidents.
\end{history}

\begin{concept}{name}
A blue box defines an idea you will need again. It always ends the same way:

\textbf{What it fixes:} the problem that made the idea worth having.

\textbf{What it costs:} what you give up by using it.
\end{concept}

\begin{question}{the question you were embarrassed to ask}
A teal box asks the question that occurs to anyone reading the paragraph above
it and that nobody asks out loud, and then answers it.
\end{question}

\begin{reflection}{the thing that should bother you}
A purple box stops on a statement that sounds absurd, or that the text just said
in passing as if it were obvious, and works out what it means and why it is
true.
\end{reflection}

\begin{etymology}
A gray box says where a word comes from. The origin of a word is not its
definition, but it often shows what the people who coined it were looking at.
\end{etymology}

The text between the boxes is the argument that connects them.

\section*{Notation}

Quantum states are written in the notation Paul Dirac introduced in 1939: $\ket{\psi}$
for a state, $\bra{\phi}$ for its dual, $\braket{\phi|\psi}$ for the inner product of
the two. You do not need to know it before starting. Chapter~\ref{ch:quantum} says
where it came from, and Chapters~\ref{ch:hilbert} and~\ref{ch:state} say what it
means.

\section*{Numbers and figures}

There is almost no code in this book, but every number in it was computed before it
went in. When a figure shows data, such as random points or a curve, the data come
from a real run with a fixed seed, and the script that produced them sits next to the
figure in the book's sources.

\section*{Exact facts, firm opinions}

This is not a scientific book, and it does not pretend to be one.

The history is checked. When the book says who did something and when, that is
what happened. Hardware numbers, error rates and resource estimates change every
few months, so every one of them comes with the date it was true.

The opinions are another matter. When a press release announces a qubit nobody has
shown to exist, the book calls it marketing. When a quantum machine learning paper does
not say how the data get into the qubits, the book says the speedup is not there. When
an analogy misleads, like the coin that is heads and tails at once, the book says so,
even when an earlier chapter used it. Sometimes that will be unfair or imprecise, and
we accept that.

A reader who leaves with a position can argue with it, refine it, or throw it
away. A reader who leaves with ``it depends'' has nothing to hold. Between being
precise and leaving you with a position, this book prefers to give you a
position.
